\documentclass[a4paper,12pt]{article}
\usepackage[T2A]{fontenc}
\usepackage[utf8]{inputenc}
\usepackage[russian]{babel}
\usepackage{amsmath,amssymb,mathtools}
\usepackage{helvet}
\renewcommand{\familydefault}{\sfdefault}
\usepackage{icomma}
\usepackage{geometry}
\geometry{left=22mm,right=22mm,top=22mm,bottom=22mm}
\usepackage{microtype}
\usepackage{tikz}
\usetikzlibrary{arrows.meta,positioning,shapes.geometric}
\usepackage{booktabs,tabularx,array}
\usepackage{enumitem}
\usepackage{hyperref}
\usepackage{parskip}
\usepackage{needspace}
\usepackage{fancyhdr}
\usepackage{setspace}
\usepackage{xcolor}
\definecolor{accent}{HTML}{008080}
\definecolor{darkaccent}{HTML}{00545A}
\definecolor{textgray}{HTML}{30383A}
\hypersetup{hidelinks,pdftitle={Остатки при делении и чётность чисел — чек-лист},pdfauthor={Лёвин Артём Александрович}}
\setlength{\parindent}{0pt}
\setlength{\parskip}{5pt}
\setlength{\headheight}{24pt}
\onehalfspacing
\setcounter{secnumdepth}{0}
\setlist[itemize]{leftmargin=6mm,itemsep=1pt,topsep=2pt}
\setlist[enumerate]{leftmargin=8mm,itemsep=11pt,topsep=4pt}
\pagestyle{fancy}
\fancyhf{}
\fancyhead[L]{\footnotesize\color{textgray}Остатки и чётность чисел}
\fancyhead[R]{\footnotesize\color{textgray}Ксения Клыкова}
\fancyfoot[C]{\small\thepage}
\renewcommand{\headrulewidth}{0.3pt}
\renewcommand{\headrule}{\hbox to\headwidth{\color{accent}\leaders\hrule height \headrulewidth\hfill}}
\newcommand{\sect}[1]{\Needspace{9\baselineskip}\vspace{5mm}{\Large\color{darkaccent}#1}\par\vspace{1mm}{\color{accent}\rule{\linewidth}{0.5pt}}\vspace{1mm}}
\newcommand{\idea}[1]{\textcolor{darkaccent}{\textbf{#1}}}
\newcommand{\note}[1]{\textcolor{darkaccent}{\textbf{Обратите внимание.}} #1}
\begin{document}
\thispagestyle{empty}
\begin{center}
\vspace*{20mm}
{\large\color{darkaccent}МАТЕМАТИКА}\par
\vspace{15mm}
{\LARGE Чек-лист}\par
\vspace{12mm}
{\Large\color{accent}Остатки при делении\par и чётность целых чисел}\par
\vspace{12mm}
{\normalsize Введение в теорию чисел\par Подготовка к экзаменационным и олимпиадным задачам}\par
\vfill
{\normalsize Ксения Клыкова}\par
\vspace{8mm}
{\normalsize 14 ноября 2026 г.}\par
\vspace{18mm}
{\small Лёвин Артём Александрович\par эксклюзивно для Ксении Клыковой}\par
\vspace{17mm}
\begin{tikzpicture}
  \draw[accent,thick] (0,0) .. controls (1.2,0.55) and (2.9,-0.5) .. (4.2,0)
                .. controls (5.5,0.55) and (7.2,-0.5) .. (8.5,0);
\end{tikzpicture}
\end{center}
\clearpage
\sect{1. Деление с остатком и сравнения}
Для любого целого числа $a$ и натурального $m$ существует единственное представление
\[
\boxed{a=mq+r,\qquad q\in\mathbb Z,\quad 0\le r<m.}
\]
Здесь $q$ --- неполное частное, $r$ --- остаток. Например:
\[
3=2\cdot1+1,\qquad 5=2\cdot2+1.
\]
Остатки чисел $3$ и $5$ при делении на $2$ совпадают. Это обозначают так:
\[
3\equiv5\pmod{2}.
\]
\idea{Определение:} $a\equiv b\pmod m$ означает, что $a$ и $b$ дают одинаковый остаток при делении на $m$. Равносильная проверка:
\[
\boxed{a\equiv b\pmod m\quad\Longleftrightarrow\quad m\mid(a-b).}
\]
Знак $m\mid(a-b)$ читается: «$m$ делит $a-b$».

\sect{2. Периодичность остатков}
При делении на $m$ возможны только остатки $0,1,\ldots,m-1$. Они повторяются с периодом $m$.

\begin{center}
\small
\renewcommand{\arraystretch}{1.25}
\begin{tabular}{c*{10}{c}}
\toprule
Число &0&1&2&3&4&5&6&7&8&9\\
\midrule
Остаток по $2$ &0&1&0&1&0&1&0&1&0&1\\
Остаток по $3$ &0&1&2&0&1&2&0&1&2&0\\
\bottomrule
\end{tabular}
\end{center}
\idea{Пример.} Найдём остаток числа $250$ при делении на $3$:
\[
250=3\cdot83+1\quad\Longrightarrow\quad250\equiv1\pmod3.
\]
Периодичность помогает найти ответ без деления большого числа столбиком.

\sect{3. Чётные и нечётные числа}
Чётное целое число имеет вид $2k$, нечётное --- $2k+1$, где $k\in\mathbb Z$.
\[
\begin{array}{c@{\quad\Longleftrightarrow\quad}c}
 n\text{ чётно} &n\equiv0\pmod2,\\
 n\text{ нечётно} &n\equiv1\pmod2.
\end{array}
\]
Форма $2k-1$ также задаёт любое нечётное число (с другим целым значением $k$).

\idea{Пример.} Для любого $m\in\mathbb Z$:
\[
4m+3=2(2m+1)+1\quad\text{--- нечётно},
\]
\[
4m+6=2(2m+3)\quad\text{--- чётно}.
\]
\sect{4. Действия с чётными и нечётными числами}
Пусть Ч --- чётное число, Н --- нечётное. Для суммы и разности работают одинаковые правила.
\begin{center}
\renewcommand{\arraystretch}{1.3}
\begin{tabular}{>{\centering\arraybackslash}p{0.28\linewidth}>{\centering\arraybackslash}p{0.28\linewidth}>{\centering\arraybackslash}p{0.28\linewidth}}
\toprule
Действие & Результат & Причина по модулю $2$\\
\midrule
Ч $\pm$ Ч & Ч & $0\pm0\equiv0$\\
Ч $\pm$ Н & Н & $0\pm1\equiv1$\\
Н $\pm$ Н & Ч & $1\pm1\equiv0$\\
Ч $\cdot$ Ч, Ч $\cdot$ Н & Ч & есть множитель $0$\\
Н $\cdot$ Н & Н & $1\cdot1\equiv1$\\
\bottomrule
\end{tabular}
\end{center}
\idea{Строгое доказательство.} Пусть $a=2r+1$ и $b=2s+1$, где $r,s\in\mathbb Z$. Тогда
\[
a+b=2(r+s+1),\quad a-b=2(r-s),
\]
\[
ab=(2r+1)(2s+1)=2(2rs+r+s)+1.
\]
Следовательно, сумма и разность двух нечётных чисел чётны, а произведение нечётно.

\sect{5. Арифметика остатков}
Если $a\equiv r\pmod m$, $b\equiv s\pmod m$, то
\[
\begin{aligned}
a+b&\equiv r+s\pmod m,\\
a-b&\equiv r-s\pmod m,\\
ab&\equiv rs\pmod m,\\
a^n&\equiv r^n\pmod m\qquad(n\in\mathbb N).
\end{aligned}
\]
После действия остаток приводят к одному из чисел $0,1,\ldots,m-1$.

\idea{Пример.} Определим чётность числа
\[
2026^{17}+35^{24}-18^3.
\]
Так как $2026\equiv0$, $35\equiv1$, $18\equiv0\pmod2$, получаем:
\[
0^{17}+1^{24}-0^3\equiv1\pmod2.
\]
\idea{Ответ:} число нечётное.

\note{Прибавление или вычитание чётного числа сохраняет чётность. Прибавление или вычитание нечётного меняет её.}

\sect{6. Два и три последовательных целых числа}
Среди $n$ и $n+1$ обязательно есть чётное число, поэтому
\[
\boxed{2\mid n(n+1)}\qquad(n\in\mathbb Z).
\]
Среди $n,n+1,n+2$ есть число, кратное $2$, и число, кратное $3$. Так как $2$ и $3$ взаимно просты,
\[
\boxed{6\mid n(n+1)(n+2)}.
\]
\note{Для проверки делимости на $6$ достаточно доказать делимость одновременно на $2$ и на $3$.}

\sect{7. Инвариант чётности}
Если исходное число нечётно, а каждое действие прибавляет или вычитает чётное число, все последующие результаты останутся нечётными.

\idea{Пример.} Начинаем с $7$. Разрешены операции $+4$ и $-6$ в произвольном порядке. Получить $100$ невозможно: исходное число нечётно, все изменения чётны, а $100$ чётно.

\sect{8. Сколько чётных чисел в последовательности?}
В каждых двух последовательных целых числах ровно одно чётное. В $31$ последовательном числе, начиная с чётного, будет $16$ чётных и $15$ нечётных.

\idea{Пример.} Среди чисел от $2$ до $32$ включительно --- $16$ чётных и $15$ нечётных. Их сумма нечётна, потому что нечётных слагаемых ровно $15$. Чётные слагаемые на чётность суммы не влияют.
\clearpage
\sect{9. Квадраты и остатки}
Пусть $k\in\mathbb Z$.

\idea{Квадрат чётного числа:}
\[
(2k)^2=4k^2\quad\Longrightarrow\quad(2k)^2\equiv0\pmod4.
\]
\idea{Квадрат нечётного числа:}
\[
(2k+1)^2=4k^2+4k+1=4k(k+1)+1.
\]
Так как $k(k+1)$ чётно, число $4k(k+1)$ делится на $8$. Значит,
\[
\boxed{(2k+1)^2\equiv1\pmod8.}
\]
В частности, квадрат любого целого числа при делении на $4$ даёт остаток \textbf{только $0$ или $1$}:
\[
\boxed{x^2\equiv0\ \text{или}\ 1\pmod4.}
\]
\sect{10. Применение: невозможность уравнения}
Проверим, может ли уравнение
\[
x^2+y^2=2027,\qquad x,y\in\mathbb Z
\]
иметь решения. По модулю $4$ каждое из чисел $x^2,y^2$ равно $0$ или $1$. Поэтому возможные остатки их суммы --- $0$, $1$ или $2$.

Но $2027=4\cdot506+3$, то есть $2027\equiv3\pmod4$. Такого остатка сумма двух квадратов дать не может.

\idea{Вывод:} целочисленных решений нет.

\note{Если правая часть имеет остаток $0$, $1$ или $2$, эта проверка сама по себе не доказывает существование решений: нужны дополнительные рассуждения.}

\clearpage
\thispagestyle{fancy}
\begin{center}
{\Large\color{darkaccent}Алгоритм: проверка уравнения суммой двух квадратов}\par
\vspace{4mm}
{\normalsize Проверка на невозможность целочисленных решений по модулю $4$}
\end{center}
\vspace{10mm}
\begin{center}
\begin{tikzpicture}[
 >=Latex,
 line width=0.75pt,
 every node/.style={font=\small,align=center},
 terminal/.style={ellipse,draw=darkaccent,minimum width=53mm,minimum height=12mm,inner sep=5pt},
 process/.style={rectangle,rounded corners=2pt,draw=accent,text width=95mm,minimum height=15mm,inner sep=7pt},
 decision/.style={diamond,aspect=2.25,draw=darkaccent,inner sep=3pt,text width=39mm,minimum height=18mm},
 result/.style={rectangle,rounded corners=2pt,draw=accent,text width=47mm,minimum height=22mm,inner sep=6pt},
 link/.style={-{Latex[length=2.5mm]},draw=darkaccent}]
\node[terminal] (start) {Дано $x^2+y^2=N$,\\$x,y\in\mathbb Z$};
\node[process,below=13mm of start] (sq) {Каждый квадрат по модулю $4$\\имеет остаток $0$ или $1$};
\node[process,below=13mm of sq] (sum) {Возможные остатки суммы:\\$0+0=0$, $0+1=1$, $1+1=2$};
\node[decision,below=16mm of sum] (cond) {$N\equiv3\pmod4$?};
\node[result,below left=22mm and 7mm of cond] (yes) {Да.\par\textbf{Целочисленных\\решений нет}};
\node[result,below right=22mm and 7mm of cond] (no) {Нет.\par\textbf{Признак не решает задачу.} Продолжите исследование};
\draw[link] (start) -- (sq);
\draw[link] (sq) -- (sum);
\draw[link] (sum) -- (cond);
\draw[link] (cond.west) -| node[pos=0.35,above]{да} (yes.north);
\draw[link] (cond.east) -| node[pos=0.35,above]{нет} (no.north);
\end{tikzpicture}
\end{center}
\clearpage
\sect{11. Типичные ошибки и памятка}
\begin{itemize}
\item Ошибка: считать, что любой квадрат кратен $4$. Правильно: остаток квадрата --- $0$ или $1$.
\item Ошибка: приравнивать сами числа, если совпали остатки. Сравнение $a\equiv b\pmod m$ означает лишь делимость $a-b$ на $m$.
\item Ошибка: забывать знак минус. Например, $0-1\equiv1\pmod2$.
\item Ошибка: раскрывать скобки без всех произведений: $(2r+1)(2s+1)=4rs+2r+2s+1$.
\item Ошибка: заключать, что из $N\not\equiv3\pmod4$ автоматически следует существование целочисленных решений уравнения $x^2+y^2=N$.
\end{itemize}

\idea{Маршрут решения:} определите модуль $\to$ замените числа остатками $\to$ выполните допустимые действия $\to$ найдите итоговый остаток $\to$ сформулируйте вывод о делимости или чётности.

\clearpage
\sect{Тренировочный блок}
\textit{Путь А. Десять заданий по возрастанию сложности. Решайте самостоятельно; ответы помещены после блока.}
\begin{enumerate}[label=\textbf{\arabic*.}]
\item Заполните остатки при делении чисел $0,1,2,3,4,5,6,7,8$ на $2$ и на $3$.
\item Найдите остатки: а) числа $251$ при делении на $3$; б) числа $2031$ при делении на $4$.
\item Определите чётность выражений $4m+3$ и $4m+6$ при любом $m\in\mathbb Z$. Обоснуйте ответы формулами.
\item Докажите, что произведение трёх последовательных целых чисел всегда делится на $6$.
\item Пусть $a$ и $b$ --- нечётные целые числа. Докажите, что $a+b$ и $a-b$ чётны, а $ab$ нечётно.
\item Не вычисляя степени, определите чётность числа $2028^{17}+41^{24}-16^9$.
\item Начинают с числа $9$ и могут сколько угодно раз прибавлять $4$ или вычитать $6$. Можно ли получить число $100$? Обоснуйте.
\item Пусть $a$ нечётно, а $b$ чётно. Определите чётность $(a+b)(a-b)$ и докажите утверждение.
\item Рассмотрите $31$ последовательное целое число, начиная с $4$. Сколько среди них чётных и нечётных? Определите чётность их суммы.
\item Докажите, что уравнение $x^2+y^2=2031$ не имеет решений в целых числах.
\end{enumerate}
\clearpage
\sect{Ответы к тренировочному блоку}
\renewcommand{\arraystretch}{1.45}
\begin{tabularx}{\linewidth}{>{\centering\arraybackslash}p{12mm}X}
\toprule
№ & Ответ / ключевая проверка\\
\midrule
1 & По $2$: $0,1,0,1,0,1,0,1,0$; по $3$: $0,1,2,0,1,2,0,1,2$.\\
2 & а) $2$, так как $251=3\cdot83+2$; б) $3$, так как $2031=4\cdot507+3$.\\
3 & $4m+3$ --- нечётно; $4m+6$ --- чётно.\\
4 & Всегда делится на $6$: в тройке есть число, кратное $2$, и число, кратное $3$.\\
5 & $a+b$ --- чётно; $a-b$ --- чётно; $ab$ --- нечётно.\\
6 & Нечётное: по модулю $2$ получаем $0+1-0\equiv1$.\\
7 & Нет: из нечётного числа прибавлением/вычитанием чётных нельзя получить чётное.\\
8 & Нечётно: оба множителя нечётны; их произведение нечётно.\\
9 & $16$ чётных и $15$ нечётных; сумма нечётна.\\
10 & Решений нет: квадраты по модулю $4$ дают $0$ или $1$, их сумма --- $0$, $1$ или $2$; $2031\equiv3\pmod4$.\\
\bottomrule
\end{tabularx}
\vfill
\begin{center}
{\small\color{textgray}Лёвин Артём Александрович --- эксклюзивно для Ксении Клыковой}
\end{center}
\end{document}
